\section{Prompt Construction}\label{app_sec:prompt_construction}

In this section, we provide more details on the prompt construction process.
The prompt used in our experiments is composed of two parts:
\begin{enumerate}
    \item \textit{instruction}: a short sentence that describes the task.
    \item \textit{demonstration examples}: a set of examples that demonstrate the expected behavior of the model.
\end{enumerate}
%
We observe that the performance of the model is sensitive to the wording of the instruction and the format of the demonstration examples.
To find a good instruction and demonstration examples, we first manually construct a set of instructions and demonstration formats.
Then, we randomly sample a few demonstration examples from the training set and manually check whether the model can solve the task with the given instruction and demonstration examples.
This process helps us find a good instruction and demonstration format, thought probably not the best.
Since our goal is not to find the best instruction and demonstration format, we didn't spend too much time on this process.
%
Below, we provide more details on the prompt construction process for each task.

\subsection{Information Extraction}\label{app_subsec:information-extraction}
The instruction used is \textbf{\texttt{Extract the triples in subject-collapsed format from texts below}}.
%
The demonstration examples are in the following format:
\textbf{\texttt{[input] -> [output]}}, where the input is a text and the output is a set of triples in subject-collapsed format.
A concrete example would be: \\
\textbf{\texttt{Vettaikaaran (2009 film) was originally written in the Tamil language, with B. Babusivan as the screenwriter. -> [s]
    Vettaikaaran\_(2009\_film) [r] original language of film or TV show [o] Tamil\_language [r] screenwriter [o] B.\_Babusivan [e]}}


\subsection{Entity Disambiguation}\label{app_subsec:entity-disambiguation}

We show two prompt construction methods for entity disambiguation.

\xhdr{Prompt construction A}
The instruction  is \textbf{\texttt{Disambiguate the entity surrounded by [START\_ENT] and [END\_ENT] by giving the correct entity name in the knowledge base}}.
%
We randomly select some demonstration examples from the training set.
We use the following format to represent the demonstration examples as a string:
\textbf{\texttt{[input] -> [mention] [ [target entity] ]}}, where the input is a text, the mention is the entity mention surrounded by [START\_ENT] and [END\_ENT], and the target entity is the entity name in the knowledge base.
%
A concrete example of demo example representation would be:
\textbf{\texttt{Eu rejects [START\_ENT] German [END\_ENT] call to boycott British lamb Peter Blackburn Brussels 1996 08\" -> German [ Germany ]}}
%
The final prompt is a concatenation of the instruction and the demonstration examples.
A full prompt with 2 demo examples would be: \\
\textbf{\texttt{Disambiguate the entity surrounded by [START\_ENT] and [END\_ENT] by giving the correct entity name in the knowledge base: "Eu rejects [START\_ENT] German [END\_ENT] call to boycott British lamb Peter Blackburn Brussels 1996 08" -> German [ Germany ]; 16 other items that were put up for auction by [START\_ENT] Hendrix [END\_ENT] s former girlfriend Kathy Etchingham -> Hendrix [ Jimi Hendrix ]; lead with a well struck header in the seventh minute [START\_ENT] Japan [END\_ENT] then laid siege to the Syrian penalty area for most -> "}}

\xhdr{Prompt construction B}
The instruction is:
\textbf{\texttt{Disambiguate the entity surrounded by [START\_ENT] and [END\_ENT] by giving the canonical entity name in the knowledge base:}}.
%
The demonstration examples are in the following format:
\textbf{\texttt{[input] -> [mention] [ [target entity] ]}}, where the input is a text, the mention is the entity mention surrounded by [START\_ENT] and [END\_ENT], and the target entity is the canonical entity name in the knowledge base.
%
A concrete example of demo example would be:
\textbf{\texttt{"Eu rejects [START\_ENT] German [END\_ENT] call to boycott British lamb Peter Blackburn Brussels 1996 08 -> German : Canonical form [ Germany ]}}
%
A full prompt with 2 demo examples would be:
\textbf{\texttt{Disambiguate the entity surrounded by [START\_ENT] and [END\_ENT] by giving the canonical entity name in the knowledge base: "Eu rejects [START\_ENT] German [END\_ENT] call to boycott British lamb Peter Blackburn Brussels 1996 08 -> German : Canonical form [ Germany ] 16 other items that were put up for auction by [START\_ENT] Hendrix [END\_ENT] s former girlfriend Kathy Etchingham -> Hendrix : Canonical form [ Jimi Hendrix ] lead with a well struck header in the seventh minute [START\_ENT] Japan [END\_ENT] then laid siege to the Syrian penalty area for most -> "}}


\paragraph{Prompt used in each dataset}

We compare the performance of the two prompts in each dataset over the validation set.
We select the prompt that performs the best in each dataset.
The prompt used in each dataset is shown in Table~\ref{tab:prompt}.

\begin{table}[h!]
    \centering
    \begin{tabular}{l|l}
        \toprule
        Dataset & Prompt \\
        \midrule
        \textsc{AIDA} & Prompt 1 \\
        \textsc{MSNBC} & Prompt 2 \\
        \textsc{AQUAINT} & Prompt 2 \\
        \textsc{ACE2004} & Prompt 2 \\
        \textsc{CWEB} & Prompt 2 \\
        \textsc{WIKI} & Prompt 1 \\
        \bottomrule
    \end{tabular}
    \caption{Prompt used in each dataset}
    \label{tab:prompt}
\end{table}

\subsection{Constiutency Parsing}
The instruction is:
\textbf{\texttt{Perform constituency parsing on the provided sentences in accordance with the Penn TreeBank annotation guidelines.}}.
%
The demonstration examples are in the following format:
\textbf{\texttt{[input] -> [output]}}, where the input is a text and the output is the flattened constituency parse tree.
A concrete example would be: \\
\textbf{\texttt{The dollar weakened against most other major currencies -> [ ( S ( NP-SBJ ( DT The ) ( NN dollar ) ) ( VP ( VBD weakened ) ( PP ( IN against ) ( NP ( RBS most ) ( JJ other ) ( JJ major ) ( NNS currencies ) ) ) ) ) ]}},
where the corresponding parse tree is shown in \Listingref{lst:parse_tree}.

\lstinputlisting[caption={parse tree with hierarchical indentation}, basicstyle=\footnotesize,label={lst:parse_tree}]{code/parse_tree.txt}


One variant of the prompt is to provide a reasoning chain with the demonstration examples.
A concrete example would be: \\
\textbf{\texttt{"The constituency parse tree is: ( S ( NP-SBJ ( NN Bond ) ( NNS prices ) ) ( VP ( VBD were ) ( ADJP-PRD ( RB barely ) ( JJR higher ) ) ) )\\"
                  "S: This stands for Sentence, the top-level structure in the parse tree."\\
                  "NP-SBJ: This is the subject noun phrase of the sentence, which is 'Bond prices'."\\
                  "NN: This stands for Noun, Singular or Mass. In this case, the word 'Bond' falls into this category."\\
                  "NNS: This stands for Noun, Plural. The word 'prices' is an example of this."\\
                  "VP: This stands for Verb Phrase, which in this case is 'were barely higher'."\\
                  "VBD: This stands for Verb, Past Tense. The word 'were' falls into this category."\\
                  "ADJP-PRD: This stands for Adjective Phrase, used as a predicate. The phrase 'barely higher' is an example of this."\\
                  "RB: This stands for Adverb. The word 'barely' falls into this category."\\
                  "JJR: This stands for Adjective, Comparative. The word 'higher' falls into this category.'"}}